\documentclass[11pt]{article}
\usepackage[margin=1in]{geometry}
\usepackage[T1]{fontenc}
\usepackage{lmodern,microtype,graphicx,booktabs,longtable,array,amsmath}
\usepackage[round,authoryear]{natbib}
\usepackage[hidelinks]{hyperref}
\usepackage{placeins}
\graphicspath{{build/}}
\input{build/numbers.tex}
\title{Geographic variation in local-government reservation assignment}
\author{Gaurav Sood}
\date{Research note --- \today}
\begin{document}
\maketitle
\begin{abstract}
Statewide descriptions of reservation assignment leave open how assignment patterns vary across administrative units. We examine women's reservation transitions in Rajasthan and Uttar Pradesh and disagreements with supplied assignment predictions for an Uttar Pradesh election. In Rajasthan's latest observed transition, district dispersion is \RajDistrictObserved\ percentage points, compared with \RajDistrictBenchmark\ under a homogeneous-transition benchmark. High statewide agreement with supplied predictions also coexists with geographic concentration of disagreements. These comparisons concern observed assignment patterns in the available data. The benchmarks do not reconstruct the complete legal allocation mechanism, and geographic concentration does not identify bureaucratic discretion. The evidence motivates a documentary investigation of where a stated rule yields different local outcomes and where its implementation differs.
\end{abstract}

\section{The finding and its interpretation}
Reservation patterns vary across districts and blocks within states. The question is how much of that variation exceeds what a specified reference process would produce, and which part remains after accounting for the actual allocation procedure. A statewide average, balance test, or strong assignment discontinuity does not directly answer that question.

We present two complementary descriptive exercises. The first follows women's reservation status across election cycles and compares geographic dispersion in transition contrasts with a common-probability benchmark. The second examines geographic concentration of disagreements between observed Uttar Pradesh reservations and predictions supplied in a published replication archive \citep{cdm2025}. Their measurements differ: a transition contrast describes repetition across cycles, whereas a prediction disagreement compares two assignments for the same election. We do not combine them into a single compliance index.

The sharpest transition result is in Rajasthan. In the 2015--2020 transition, the standard deviation across district contrasts is \RajDistrictObserved\ percentage points, compared with a benchmark mean of \RajDistrictBenchmark. Across blocks the corresponding figures are \RajBlockObserved\ and \RajBlockBenchmark. Uttar Pradesh's latest district spread exceeds its reference range, while its block spread lies within that range. The supplied-prediction exercise establishes that low statewide disagreement can nevertheless be geographically concentrated.

The contribution is to measure this dispersion and examine its sensitivity, not to establish that earlier studies overlooked assignment. The evidence appendix distinguishes studies \citep{chauchard2014,dunning2013,das2023,cdm2025,weaver2024} that describe procedures, assess balance, reconstruct assignment, or examine implementation heterogeneity. It also distinguishes their populations and election cycles. Our results do not invalidate their outcome estimates: that would require examining their identifying comparisons, the location of discrepancies relative to those comparisons, and the relevant outcome assumptions. Reexamining identifying claims across periods is also central to \citet{green2023}; the related reservation study by \citet{soodquota} motivates attention to the assignment process underlying comparisons across cycles.

\section{Institutional expectations}
For the offices examined here, the relevant institutional hierarchy runs from the Indian Constitution through state legislation and rules to applicable administrative orders \citep{constitution,rajasthanact,rajasthanrules,upact,uprules,up2021}. Constitutional requirements and delegation to states do not themselves establish authority for districts to adopt different procedures. Equally, assigning an officer responsibility for implementing a procedure does not imply unrestricted discretion over its content. The appendix records the legal provisions and the limits of their historical applicability.

A common procedure need not imply identical unconditional transition rates. The eligible pool, quota, reservation history, caste-category allocation, and lottery realization can differ across places. These are empirical explanations to test, not reasons to dismiss geographic heterogeneity. An adequate explanation must show that the same applicable procedure, supplied with the correct local inputs, generates the relevant pattern.

Conversely, a procedural departure requires comparison with the operative rule and the information it required. A prediction error against a researcher's reconstruction is not automatically such a departure. Where several historical reconstructions remain plausible, the appropriate output is a range of compatible assignments or a classification as unresolved. A contemporaneous official list, a modern population aggregate, and an author's predicted assignment are distinct objects.

\section{Data and measurement}
The transition exercise reads checksum-pinned Rajasthan and Uttar Pradesh master panels from the upstream \texttt{local\_elections} repository. The analysis identifies GP-head observations and links reservation status across configured general-election cycles. Geographic identities and matching decisions are part of the measurement process. Early Rajasthan coverage is particularly incomplete. Consequently, the estimands concern the linked observations that enter each comparison, not an established census of historically unchanged GPs.

The second exercise uses the Chaturvedi--Das--Mahajan replication data. The upstream study repository supplies this input. There are \AuditN\ records with observed assignments and the required geography in the audit sample. The primary women's prediction is the supplied instrument; a separate supplied prediction is retained as a sensitivity. SC and ST assignments are examined separately. Missing observations retain their exclusion status rather than being recoded as unreserved.

The Census linkage extension uses upstream village--GP mappings and SHRUG population data. Weaver's March 2025 source and prepared wide panel reside in \texttt{local\_elections\_up}; they are not an unpublished dependency of this repository. Modern GP boundaries do not establish historical boundaries. The linkage therefore informs future reconstruction and is not used to certify historical rank-rule compliance in this note.

\section{Estimands and reference distributions}
For geographic unit $g$ and transition $t$, define
\[
\Delta_{gt}=\Pr(W_t=1\mid W_{t-1}=1,g)-\Pr(W_t=1\mid W_{t-1}=0,g),
\]
where $W$ indicates reservation of the office for a woman. We estimate each probability by its observed linked-sample proportion. A negative contrast describes less immediate repetition; it does not, by itself, establish lawful rotation. Units without observations in both lag groups cannot identify this contrast and are excluded from its dispersion calculation.

The primary geographic spread is the equal-unit standard deviation of these contrasts. The benchmark probabilities are estimated from all linked seats in the transition, whereas dispersion requires both lag groups within a geographic unit. The audit evaluates this support difference. District contrasts are computed directly from their GP observations. Block contrasts are calculated within blocks nested in districts. Seat-weighted summaries answer a different question and enter the sensitivity analysis. The observed dispersion includes variation associated with finite group sizes.

The transition benchmark fixes each included unit's two lag-group denominators and independently draws current reservations using statewide probabilities conditional on prior reservation. It estimates the spread produced by that homogeneous Bernoulli process. It does not preserve local quotas, reconstruct eligibility, or simulate a sequence of legally constrained lotteries. Its simulated range is a reference distribution conditional on fitted probabilities, not a confidence interval for a latent geographic variance.

For the supplied-prediction exercise, define $M_i=1$ when observed and predicted reservation disagree. We report its statewide mean, geographic rates, and their dispersion. The reference distribution redistributes observed disagreement indicators within caste/prediction strata, preserving stratum totals and geographic sample sizes. This addresses concentration relative to exchangeable disagreements. It does not preserve every feature of a feasible joint assignment and is not a randomization test of the legal allocation mechanism.

We also decompose observed block-rate variance into between-district and within-district components using the same weights throughout the identity. The between component uses district averages of block rates under those weights; it is not interchangeable with the separately calculated dispersion of direct district rates. The resulting share includes finite-group variation and must not be interpreted as the fraction of underlying variation caused by district administrations.

\section{Geographic results}
Figure~\ref{fig:transitions} compares observed transition dispersion with the homogeneous-process reference. Rajasthan's latest transition has substantial excess dispersion at both administrative levels. Its earliest observed transition lies within the reference range at both levels; excess dispersion appears in the two later transitions. This timing motivates investigation of period-specific inputs and procedures rather than a presumption of fixed district behavior. Uttar Pradesh's latest district spread is \UpDistrictObserved\ percentage points against \UpDistrictBenchmark\ in the reference; its block spread is \UpBlockObserved\ against \UpBlockBenchmark. The comparison therefore does not support a blanket statement that every state and administrative level exhibits equally large departures.

\begin{figure}[htbp]
\centering\includegraphics[width=\linewidth]{transition_spread.pdf}
\caption{Dispersion of women's reservation-transition contrasts. The universe is linked GP-head observations with both lag groups represented within each included unit. Points show observed spread and the homogeneous-process benchmark mean; intervals show the middle 95\% of simulated benchmark spreads. Units are percentage points. The simulation preserves lag-group sizes, not the complete legal assignment mechanism.}\label{fig:transitions}
\end{figure}

Statewide disagreement with the primary women's prediction is \WomenMismatchRate\ percent. The SC and ST disagreement rates are \ScMismatchRate\ and \StMismatchRate\ percent. These rates use all available observations, including seats both sources classify as unreserved. For a rare category, high overall agreement can consequently coexist with appreciable disagreement among seats classified as reserved by either source. The generated audit table reports that alternative denominator as well.

Figure~\ref{fig:assignment} asks whether the observed disagreements are geographically concentrated relative to the stratified reference. Women's and SC disagreements show greater dispersion than that reference. For SC, \ScBalancedBlocks\ of the \ScAffectedBlocks\ blocks with disagreements have equal numbers of false-positive and false-negative assignments. This pattern is compatible with within-pool substitutions that preserve the total quota, although it does not establish their cause. ST has few disagreements and is weakly informative about geographic implementation. The primary women's instrument and alternative prediction remain separate specifications; neither is treated as an independently verified legal assignment.

\begin{figure}[htbp]
\centering\includegraphics[width=\linewidth]{assignment_spread.pdf}
\caption{Geographic spread of supplied-prediction disagreements in Uttar Pradesh. Observed spread is compared with the mean and middle 95\% range from redistributing disagreement labels within the specified strata. Geographic units receive equal weight. The reference conditions on the number of disagreements; it does not estimate their cause or simulate lawful joint caste allocation.}\label{fig:assignment}
\end{figure}

Figure~\ref{fig:distribution} makes the denominators visible. An extreme rate from a small administrative unit should not carry the same evidentiary weight as a comparable rate measured across many seats. The figure displays the observed quantities without attaching model-based confidence intervals to full-source descriptive counts. Selection into the available source universe remains a separate uncertainty.

\begin{figure}[htbp]
\centering\includegraphics[width=\linewidth]{geographic_distribution.pdf}
\caption{Prediction disagreements and geographic sample sizes. Each point represents an observed district or block. The horizontal axis reports its available GP denominator on a logarithmic scale; the vertical axis reports its disagreement rate. Facets distinguish prediction outcomes and administrative levels. These are descriptive rates, not estimated official propensities.}\label{fig:distribution}
\end{figure}
\FloatBarrier

\section{Audit and sensitivity}
The audit reconstructs the headline quantities independently and checks the distinction between the units observed and the units summarized. Its sensitivity analyses examine minimum denominators, equal-unit and seat weights, influential districts, and alternative predictions. Removing small groups changes the target population; such results must be read with their retained counts rather than presented as an automatic improvement over the primary analysis.

The reference processes impose stronger assumptions than arithmetic correctness. Independence in the transition simulation excludes within-pool quota dependence. Exchangeability in the disagreement exercise excludes geographically correlated prediction or measurement error beyond the chosen strata. Errors in a single population input or rotation baseline can change several predictions within the same allocation pool. Fixed quotas can also pair an extra reservation with an omitted reservation elsewhere in that pool. Thus geographic concentration can arise from correlated reconstruction errors even without a departure by an official. Departures from these references can therefore reflect several mechanisms. The supplementary audit reports which numerical identities were verified and which assumptions cannot be established with current inputs.

\IfFileExists{build/audit_summary.tex}{\input{build/audit_summary.tex}}{}

\section{What would explain the variation?}
The immediate finding is geographic heterogeneity unexplained by the reference processes used here. Attributing it to bureaucratic discretion requires additional evidence. A useful investigation would reconstruct complete allocation pools using contemporaneous rosters, population inputs, reservation histories, draft lists, objections, and final signed orders. Sampling high- and low-disagreement areas together with randomly selected areas would permit both mechanism discovery and an assessment of prevalence.

The investigation should distinguish incorrect research inputs, documented exceptions, different administrative interpretations, clerical or capacity failures, and deliberate choices. A district that consistently uses another rotation baseline presents a different mechanism from isolated targeted seat substitutions. Correct quota totals with different GP selections would locate the discrepancy in selection; incorrect totals would also implicate quota construction or definitions.

Political intervention is a further hypothesis. Evidence would need to show who requested or authorized changes and which actors stood to benefit before assignment. Winning candidates are downstream of reservation and cannot alone establish why an assignment was made. Likewise, changes following an officer's transfer are descriptive unless the selection and timing of transfers support a causal comparison.

The legal question remains specific: which applicable provision permits the observed departure, if any? Geographic location alone supplies no justification. The statistical question is equally specific: which documented inputs and dependence structure can account for the observed distribution? Both questions can be pursued without treating every unexplained residual as discretion or every conceivable exception as an explanation.

\clearpage
\bibliographystyle{plainnat}
\bibliography{references}
\clearpage\appendix
\section{Generated numerical exhibits}
\subsection{Transition dispersion}
\input{build/transition_table.tex}
\subsection{Supplied-prediction disagreements}
\input{build/assignment_table.tex}
\subsection{Geographic disagreement benchmarks}
\input{build/assignment_benchmark_table.tex}

Reference ranges are the middle 95\% of the stratified disagreement permutations; all spreads are in percentage points and use equal geographic-unit weights. They are not confidence intervals.
\subsection{Population-linkage coverage}
\input{build/coverage_table.tex}
\IfFileExists{build/audit_table.tex}{\subsection{Audit sensitivities}\resizebox{\linewidth}{!}{\input{build/audit_table.tex}}

Each row uses the latest transition for its state. Minimum is the required GP count in each prior-reservation group. Units and GPs describe retained support; SD is the equal-unit standard deviation in percentage points. LODO low and high are the smallest and largest spreads after leaving out one district; these are influence ranges, not confidence intervals.}{}
\clearpage
\section{Assignment claims and the evidence behind them}
\label{app:evidence}

A description of a reservation rule, a balance test, an assignment first stage,
and a reconstruction of every seat's eligibility answer different questions.
This appendix records those distinctions. It covers the principal studies
identified through the project's Rajasthan/UP literature ledger, the
\texttt{quota\_spending} literature collection, publisher and author records,
and reference searches for reservation, randomization, rotation, and allocation.
The search includes foundational women's-reservation studies and contrasting
rank-based procedures outside Rajasthan and UP. It is a documented starting
set, not a claim that every published assignment statement has been located.
Locations below refer to the specified version; working-paper and journal
pagination are not interchangeable.

\subsection{Foundational lottery and rotation designs}

\paragraph{Chattopadhyay and Duflo.}
\citet{chattopadhyay2004} describe women's reservation as a randomized policy
experiment. The published abstract states that head positions ``have been
randomly reserved for a woman.'' The institutional discussion and balance
comparisons give the assignment claim empirical content. The relevant
comparison spans West Bengal and Rajasthan, but the field sampling units differ:
the Rajasthan hamlet sample should not be treated as an enumeration of GP
allocation pools. The conclusion needed here is narrower than a claim about
all subsequent elections: their account concerns the institutional arrangements
and samples they studied. An audit of later cycles must recover the later
eligible pools and rotation histories.

\paragraph{Beaman and coauthors.}
The abstract of \citet{beaman2009} states that they ``exploit random assignment
of gender quotas for leadership positions on Indian village councils.'' The
study uses prior exposure to women leaders to examine subsequent electoral
outcomes and attitudes. \citet{beaman2012} use the same broad policy setting to
study aspirations and attainment. Both make assignment and exposure history
central to interpretation. Neither publication's abstract alone establishes a
statewide compliance rate or the distribution of deviations across districts.
For the latter question, one needs complete administrative rosters and a
benchmark preserving the actual assignment constraints, rather than treating
an outcome effect or balance result as a compliance statistic.

\paragraph{Bhavnani.}
\citet{bhavnani2009} describe ``randomly chosen seats in local legislatures''
in the published abstract. The application concerns Mumbai municipal elections,
not Rajasthan or UP GP-head offices. It is relevant to the persistence argument
and to the use of quotas as natural experiments, but its assignment account
cannot validate a rural state's allocation mechanism. Municipal wards and GP
headships remain separate institutional universes.

\subsection{Rajasthan caste and UP assignment studies}

\paragraph{Chauchard.}
\citet{chauchard2014} describes rotation through ranked SC population shares
in Rajasthan (author PDF, pp.~7--8 and Table~1). The design uses
matched GPs near local allocation thresholds. The project's prior reading of
online Appendix~B, pp.~5--11, records checks of timing disagreements and
sensitivity to Census vintages, household versus person counts, rounding, and
boundary changes. The supplementary PDF was not successfully re-fetched during
this drafting pass; these appendix details retain that access limitation.
The main paper and the existing ledger support a comparison of assignment
validation strategies, not the assertion that Chauchard ignored implementation
imperfections. Modern Census shares on modern GP boundaries cannot by themselves
refute an election-year allocation account.

\paragraph{Dunning and Nilekani.}
\citet{dunning2013} use state-specific allocation thresholds in Karnataka,
Rajasthan, and Bihar. The December 2012 author version, pp.~20--22 and Table~2,
reports balance checks and discusses variation in implementation across
subdistricts. The account uses population totals in passages concerning
Rajasthan/Bihar, whereas Chauchard describes population shares. This distinction
must be resolved against operative orders and administrative population lists;
it must not be silently standardized by the researcher. Their checks address
comparability in the study samples. The present geographic spread statistics
address another quantity and do not overturn those checks merely by differing
from a homogeneous-process benchmark.

\paragraph{Das, Mukhopadhyay, and Saroy.}
\citet{das2023}, identification discussion and Tables~1--2 and D1, study
Rajasthan's OBC reservation after SC/ST assignments. Their lottery account
applies to the remaining eligible GPs. They examine covariate balance with block
fixed effects and additional composition checks. This is a substantial
administrative sample, so describing all previous studies as small-sample
exercises would be inaccurate. Balance conditional on blocks and geographic
variation in assignment residuals are different estimands. A reconstructed
benchmark must retain the exclusion of seats already allocated to other
categories.

\paragraph{Chaturvedi, Das, and Mahajan.}
\citet{cdm2025}, pp.~760--763 and Figure~1, describe block-specific caste and
women's assignment in UP's 2015 election using SEC population information.
They report almost 98\% prediction of female assignment and examine the
assignment discontinuity and covariate balance. This note audits the geographic
distribution of disagreements with their supplied prediction fields. Those
fields are valuable replication inputs, but their presence does not establish
that this project has independently reconstructed the legal rule. The statewide
assignment universe, near-cutoff outcome sample, and complete-case residual
audit have different denominators.

\subsection{Evidence that ranking and implementation deserve separate treatment}

\paragraph{Gajwani and Zhang.}
\citet{gajwani2015}, published pp.~239--240, explicitly state: ``the assignment
of reservations in Tamil Nadu is not random.'' Their account orders GPs using
caste and female population ratios within blocks, allocates women within caste
categories, and retains reservations across two terms before rotating. Their
main comparison and RD sensitivity should therefore not be described as
unconditional lottery assignment. This is direct evidence that a national
reservation mandate does not imply one nationwide randomization algorithm.
It says nothing by itself about the authorization of different procedures
within Rajasthan or within UP.

\paragraph{Pathak and Macours.}
\citet{pathak2017}, pp.~746--749, describe Andhra Pradesh rotation and ranking
by the female population ratio. Their outcome design compares the timing of
children's exposure to reserved leadership across cohorts. Assignment histories,
not simply current reservation indicators, define the comparison. A legal
rotation schedule can induce dependence across cycles; rejecting independence
is consequently insufficient to establish a departure from the prescribed
procedure. The DOI was verified against Crossref as
\texttt{10.1086/692114}; this corrects a conflicting DOI in an upstream
bibliography without altering the upstream repository.

\paragraph{Narasimhan and Weaver.}
\citet{weaver2024} study GP delimitation, not caste or women's reservation.
The November 2023 working-paper Appendix~OA.2, pp.~80--81, examines district
first stages and constraints associated with neighboring villages' population
viability and contiguity. The working paper also reports that district
noncompliance classifications change between delimitation rounds. This is a
precedent for investigating geographic implementation variation and for testing
concrete alternative explanations. Its particular viability constraints cannot
simply be imported into a reservation model. Likewise, classifications based on
first-stage significance depend on precision and should not become this note's
measure of latent compliance.

\subsection{Related work by the author}

\citet{soodquota}, public author version, section on rotation and its
implications and Appendix~B, studies women's electoral performance following
reservation in Rajasthan and UP. It documents transition dependence and
recognizes that conditioning on a subsequent open seat can select on an outcome
of prior assignment. The present note separates that dependence from geographic
heterogeneity and from verified legal noncompliance. Failure to reject an
independence test is not a certificate of random assignment; a negative lag
coefficient is not an independently established rotation rule.

\citet{green2023} provide a methodological comparison: extending a natural
experiment across earlier and later congresses changes the interpretation of
an apparent daughters effect. The relevance is the discipline of testing
whether a finding generalizes across periods and cohorts. That paper supplies
no evidence about Indian reservation law or administrative discretion.

\section{Legal provisions and historical applicability}
\label{app:law}

The hierarchy relevant to these two states is the Constitution of India, state
legislation, delegated rules, and applicable orders. A state-level mandate does
not automatically authorize districts to choose different allocation procedures.
The distinction between an authorized implementing officer and discretion over
the governing procedure is central to the research question.

\subsection{The constitutional mandate}

Article~243D of \citet{constitution} distinguishes seats within Panchayats
from offices of their chairpersons. Clause~(4), relevant to GP heads, assigns
the manner of chairperson reservation to state legislation and separately
requires rotation. Its final proviso states:
\begin{quote}
``the number of offices reserved under this clause shall be allotted by
rotation to different Panchayats at each level.''
\end{quote}
The clause sets a minimum women's share and population-based SC/ST reservation
requirements at the state and tier level. It does not supply every eligible
pool, ordering convention, or lottery procedure needed to simulate a particular
election. Nor does rotation, standing alone, determine whether repeated
reservation is impossible after category changes, roster changes, or exhaustion
of an eligible list. Those questions require the operative state provisions.
The cited official compilation is dated 11 November 2025; its publication date
is not the enactment date of Article~243D.

\subsection{Rajasthan: prescribed procedure and delegated execution}

Section~16 of the Rajasthan Panchayati Raj Act supplies the statutory context
for reservation of chairperson offices \citep{rajasthanact}. The India Code
compilation of the election rules describes SC/ST ranking, category overlap,
and sarpanch/pradhan allocation in Rules~7--9 \citep{rajasthanrules}.
Rule~7's account uses descending population percentages and a lower eligibility
threshold, with continuation and restarting of rotating lists. Rule~8 addresses
places appearing in both SC and ST lists. Rule~9 refers office allocation to
that procedure and assigns implementation to a government-authorized officer.
Its women's allocation provision specifies ``by draw of lots.''

These are substantive restrictions on implementation. They should be encoded
rather than replaced with a generic assertion of random assignment. The
compilation does not establish which amendments and supplementary orders applied
to every election in this panel. A compliance reconstruction still needs the
complete historical roster, prescribed quota and rounding, population vintage,
previous reservations since the applicable reset, category ordering, and the
actual authority for any exception. A district-specific procedure unsupported
by those sources is not made lawful merely because a district officer used it.

\subsection{Uttar Pradesh: historical rules remain incomplete}

The state Act and the Panchayat Raj reservation rules are distinct from rules
for municipal wards or for other Panchayat tiers \citep{upact,uprules}.
The India Code extract reviewed in the project provides definitions and quota
rounding, but does not establish a complete historical implementation of
Rule~5 for all GP-head elections. Substituting a similarly named municipal or
Kshettra/Zila Panchayat rule would change the institution under study.

The official notification of 17 March 2021 changes the definition of previous
election to refer to the 2015 general election and subsequent general elections
\citep{up2021}. Its existence makes a mechanically inherited 1995 baseline
inappropriate without further verification. The sequence of orders, amendments,
and any supersession must be matched to the final allocation used for the
election. The current supplied-prediction audit concerns 2015 and is not a
reconstruction of the 2021 procedure.

\subsection{What the legal evidence can establish now}

The legal texts establish that reservation is governed by mandates and
procedures, not that every observed assignment complies with them. The statistical
results establish departures from specified reference processes, not that those
processes equal the law. The next evidentiary step is to compare complete
administrative allocation pools with the procedure and inputs actually applicable
to them. Cases should be classified as source disagreement, a resolved historical
input difference, an authorized exception with a cited provision, or an unresolved
departure. Only the last two categories concern implementation after the rule and
inputs have been established; neither category alone identifies political motive.

\section{Source validation and remaining coverage}

Article titles, author lists, years, journals, volumes, pages and DOIs were
checked against Crossref records where available, and publisher or author
records otherwise. Primary PDFs and the existing page-specific evidence ledger
were used to distinguish assignment descriptions from validation exercises.
The bibliography retains the user-authored quota paper as an undated working
paper rather than inventing a publication year. Legal sources retain their
instrument dates and explicit compilation notes.

This pass does not certify an exhaustive review or complete election-year legal
reconstruction. The remaining search includes additional state-specific
applications identified in the upstream literature collection, unavailable
supplementary materials, and historical government and district allocation
orders. Their absence limits claims about the entire literature and about legal
compliance, rather than supplying evidence for discretion. Short quotations
above identify the precise claims under discussion; surrounding paragraphs are
paraphrases, not transcriptions of inaccessible appendices or statutes.

\end{document}
